%{\color{blue}outline for this chapter:
%\begin{itemize}
%    \item problem def and improves on prev prelim explores, transition to attention and pretrained models, large-scale experiments and near-OOD
%    \item method, technical (details - hopefully illustrations at attn level, maybe more detailed math at attn level too or sample walkthrough)
%    \begin{itemize}
%        \item maybe more detailed or broken-out figure?
%        \item attn heatmaps would be ideal if possible and they look good
%    \end{itemize}
%    \item method, theory (high-level, proof, general case, maybe toy problem?)
%    \item experiments and main results, additional results
%    \item ablations and all related exploratory analyses with discussion
%\end{itemize}
%}

\textit{[ This chapter includes parts of a paper previously published as - Alexander Politowicz, Sahisnu Mazumder, and Bing Liu. “Improving Out-of-Distribution Detection using Segmented Images and Cross-View Attention Fusion”. In Proceedings of the Winter Conference on Applications of Computer Vision (WACV 2026), 2026.]}

\section{Improving Multi-View OOD Detection for the Open-World}
\label{sec:propwork-casod}

%\sysb considers all views during learning and OOD detection through computation of all individual views, all possible cross-view combinations, and a comprehensive fused feature. Furthermore, \sysb constructs and evaluates all possible ensembles over these views, resulting in \textbf{over 100 computations per forward pass}. This brute-force approach was shown to be effective, but continues to exhibit significant shortcomings:
%\begin{itemize}
%    \item Inefficient Processing - the computational complexity of \sysb OOD detection is signifcant.
%    \item Inefficient Learning - considering all cross-view features with the view-specific and comprehensive fused features during training includes significant semantic overlap with no mechanism to leverage the repeated information, wasting learning resources.
%    \item Naive Fusion - max-pooling is an outdated and naive fusion method \cite{} that may be underperforming compared to new fusion technology.
%    \item Limited Experimental Evaluation - \sysb continues to only consider CIFAR-100 as the ID dataset. There are many other large-scale dataset with higher resolution, complex images that must be considered.
%\end{itemize}
Adapting ML to clinical open-world settings requires improving a range of factors across the ML landscape. One factor that has been thoroughly emphasized until this point, and in fact is the predominant inspiration for this dissertation, is the capacity to detect OOD samples during operation. ML models with this ability induce robustness, generalizability, safety, and trust, all factors highly-valued both in applications and in society \cite{siau2018building,hendrycks2022x,hendrycks2021unsolved,amodei2016concrete,nitsch2021out}. The primary roadblock to this progress in modern OOD detection is near-OOD data \cite{yang2022openood,winkens2020contrastive}.
Another important factor, orthogonal to the present ability of OOD detection (but just as critical), is demonstrated compatibility with modern data and computational environments. It is no longer sufficient to consider approaches that evaluate exclusively on curated datasets, require extensive training, or assume access to dedicated, abundant resources. With real-world applications in mind, ML solutions must demonstrate proficiency on challenging, large-scale, near-OOD datasets, effectively utilize existing architectures and paradigms to optimize development time, and run efficiently without access to massive computational frameworks or hardware infrastructure \cite{li2025secure,wang2025enabling,yuhas2021embedded,lu2025out}.

\gls{sysa} and \gls{sysb} make progress in the first factor - improving near-OOD detection performance - but neglect the second factor. \gls{sysa} is unstable and wasteful in learning from the provided views, necessitating more experimentation and validation. \gls{sysb} is comprehensive but slow and resource-demanding, unfit for many application requirements that operate with fast-paced, lightweight systems. Both methods require extensive training of models from scratch and show results in limited experimental settings.
This sets a precedent for several improvements that serve to propel multi-view OOD detection farther into the frame of application-oriented, open-world ML. Recent trends in NLP and CV have extolled the power of Transformer mechanisms in models \cite{dosovitskiy2020image,achiam2023gpt,wang2024mutualformer}. These architectures have been verified as effective both in OOD detection and multi-view classification \cite{sun2022out,fort2021exploring,zhao2023cddfuse} (also see Chapter \ref{diss_mvood_init} Section \ref{sec:propwork-scod}). It is natural to consider utilizing these mechanisms in multi-view OOD detection, both for their potential performance benefits and for the abundance of pre-trained Transformer models mitigating necessary training. In particular, the Attention computation inherent to Transformer models (or Transformers) gives precedence to an improved fusion technique that has already been shown to be fundamental in the success of Transformers over traditional CNNs. With experimentation on extensive new settings oriented around the large-scale ImageNet dataset \cite{deng2009imagenet}, the bridge for multi-view OOD detection towards real-world application is set.

%Without loss of generality, we represent the inputs to a CVCA module as $z_{v1}$ and $z_{v2}$. A CVCA module takes $z_{v1}$ and $z_{v2}$ and extracts queries $Q_{v1} = FF_{Q,v1}(z_{v1}) \in \mathbb{R}^{T_{v1} \times d_{k}}$ from $z_{v1} \in \mathbb{R}^{T_{v1} \times d_{v1}}$ using feed-forward layer $FF_{Q,v1}$, and keys $K_{v2} = FF_{K,v2}(z_{v2}) \in \mathbb{R}^{T_{v2} \times d_{k}}$ and values $V_{v2} = FF_{V,v2}(z_{v2}) \in \mathbb{R}^{T_{v2} \times d_{w}}$ from $z_{v2}$ using feed-forward layers $FF_{K,v2}$ and $FF_{V,v2}$. Similarly, we can obtain $Q_{v2}$, $K_{v1}$, and $V_{v1}$ using feed-forward layers $FF_{Q,v2}$, $FF_{K,v1}$, and $FF_{V,v1}$. Note that in our implementation $T_{v1} = T_{v2}$ and $d_{k} = d_{w}$. These queries, keys, and values are then used to generate the cross-attention features $z_{v2 \rightarrow v1}$ and $z_{v1 \rightarrow v2}$
Attention \cite{vaswani2017attention} deserves its own introduction, both for its existing impact and for the role it plays in improved multi-view OOD detection. Details for Attention mechanisms and component details can be found in Chapter \ref{diss_bkgd} Section \ref{sec:bkgd-attn}. Attention effectively identifies which subsections of an image to attend to with respect to the optimization objective. The delineation of ``self-attention'' refers to the fact that the same input $x$ is used to compute each component ($Q$, $K$, $V$) for attention. Thus the attention mechanism references and emphasizes components of the same input image when identifying important features. Cross-attention uniquely considers two inputs and quantifies the interactions they have with each other. Given a $Q$ that is computed from a view $x_1$ and $K$, $V$ computed from a view $x_2$, feature information in $x_2$ is referenced to identify and emphasize relevant feature information in $x_1$. This allows for the model to exploit the information in \emph{multiple inputs} to learn rich, correlated feature information across inputs. These state-of-the-art processing mechanisms enable learning of detailed, informed features conducive to classification, or more relevantly, OOD detection.

%We introduce a novel method called \textbf{C}ross-view \textbf{A}ttention of \textbf{S}egmented views for \textbf{O}OD \textbf{D}etection (\textbf{\sys}) that learns class-discriminative features using the aforementioned separated views of images and effectively integrates them for OOD detection. \sys works in two steps: given the input image and its foreground and background, (1) a \textbf{frozen} pre-trained model augmented with \textbf{trainable} parameter-efficient Adapter modules \citep{hu2021lora} is used to extract per-view features from the original image, the foreground view, and the background view of the image. (2) A novel multi-layer \emph{cross-view correlation attention} (CVCA) fusion module is then employed to intelligently and cohesively integrate the ID class-discriminative information from the three view-specific features by comprehensively attending between and fusing view features, a unique deviation from prior fusion methods \cite{li2023mvhanet}. This results in a highly discriminative fused predictive feature, improving OOD detection (see Sec. \ref{sec:exp}). As foreground-background separation is not the focus of this work, we employ an off-the-shelf fast tool to generate the foreground and background views of the image (can segment an image in 0.02 secs). 
%\item We propose intelligent cross-attention fusion of segmented foreground-background images for OOD detection. This approach has not been reported before.% Although related prior work exists, as discussed above, their settings are very different from ours. 
%\item We introduce the \textbf{first} foreground-background image OOD detection method, \sys, with a novel cross-view correlation attention (CVCA) for feature fusion. CVCA differs from prior cross-attention methods by using a \textit{hierarchy of attention operations to integrate all three view features simultaneously}.
To leverage these advantages and advance multi-view OOD detection, \textbf{C}ross-view \textbf{A}ttention of \textbf{S}egmented views for \textbf{O}OD \textbf{D}etection (\textbf{\sys}) proposes effective view feature extraction using pre-trained model backbones and novel, state-of-the-art feature fusion using hierarchical cross-view attention fusion for OOD detection. Training is streamlined using parameter-efficient Adapter modules \cite{hu2021lora} on the pre-trained backbones, eliminating significant computational requirements. Fusion is implemented using novel multi-layer \textit{\gls{cvca}} that cohesively integrates ID class-relevant discriminative feature information using view-informed attention-based mechanisms. This simultaneous fusion of view feature information enables improved learning of ID discriminative features and subsequently enhances OOD detection. With the use of Adapters and pre-trained models, these advances come at relatively little computational resource cost.

%In contrast, our approach involves explicit segmentation and cross-view interaction modeling in a multi-view learning framework that helps in a more effective learning from the isolated view representations \textit{with negligible segmentation overhead}, as shown in our Experiments. 
%\item \sys processes multiple views at a very small (trainable) parameter cost using Adapters and a frozen pre-trained model, requiring 44\% fewer trainable parameters.% It achieves OOD detection performance improvements while requiring 44\% fewer trainable parameters compared to that of fully fine-tuning the backbone model, an appealing application benefit.
%\item \sys outperforms state-of-the-art OOD detection methods by significant margins. For example, using a 100-class subset of ImageNet as the ID data, \sys improves the average AUC for near-OOD by 2.08\% and for far-OOD by 2.71\% on a relative scale. 
Following from the motivating discussion, evaluation of \gls{sys} is performed in settings with demanding, near-OOD datasets. In particular, the ImageNet-1k dataset \cite{russakovsky2015imagenet} (or subsets of it) is used as the ID dataset, containing 1000 classes with varying amounts of semantic and/or covariate overlap. Various near- and far-OOD dataset are used for testing of \gls{sys}, including the notoriously difficult NINCO \cite{bitterwolf2023or} and SSB-Hard \cite{vaze2021open} datasets, specifically designed to be similar to ImageNet-1k classes \cite{zhang2023openood}. Furthermore, rigorous ablation studies are performed along with both quantitative and qualitative analyses (including a theoretical analysis) to thoroughly substantiate the performance of \gls{sys} across a variety of relevant contexts expected in real-world applications (see Section \ref{sec:propwork-casod-disc}).
